\begin{abstract}
	% Generating unbiased scene graphs by learning from heavily imbalanced data is a vital challenge for current scene graph generation (SGG). 
	% One critical problem is to improve the performance of handling input of less frequently seen relationships in the training dataset. 
	%  bridges the gap between image data and their semantic labels to benefit a variety of downstream tasks.
% 	Scene graph generation (SGG) aims to build a structured representation of a scene using objects and pairwise relationships, which benefits many downstream tasks. However, current SGG methods usually suffer from the long-tailed distribution of training data, causing suboptimal scene graphs.
% 	To address this problem, we propose resistance training using prior bias (RTPB) for scene graph generation. Specifically, RTPB uses a prior bias, named resistance bias, to encourage the model to focus on less frequent relationships during training, thus improving the model generalizability on tail categories.
% 	In addition, to further improve scene graph generation by exploring the contextual information, we design a contextual encoding backbone network, termed as Dual Transformer (DTrans).
% 	We perform extensive experiments on a very popular benchmark, VG150, to demonstrate the effectiveness of our method for unbiased scene graph generation. Specifically, our RTPB achieves an improvement of over 10\% under the mean recall evaluation criterion, and outperforms recent state-of-the-art methods with a large margin.
	\textbf{S}cene \textbf{G}raph \textbf{G}eneration (SGG) aims to build a structured representation of a scene using objects and pairwise relationships, which benefits downstream tasks. However, current SGG methods usually suffer from sub-optimal scene graph generation because of the long-tailed distribution of training data. To address this problem, we propose Resistance Training using Prior Bias (RTPB) for the scene graph generation. Specifically, RTPB uses a distributed-based prior bias to improve models' detecting ability on less frequent relationships during training, thus improving the model generalizability on tail categories. In addition, to further explore the contextual information of objects and relationships, we design a contextual encoding backbone network, termed as Dual Transformer (DTrans). We perform extensive experiments on a very popular benchmark, VG150, to demonstrate the effectiveness of our method for the unbiased scene graph generation. In specific, our RTPB achieves an improvement of over 10\% under the mean recall when applied to current SGG methods. Furthermore, DTrans with RTPB outperforms nearly all state-of-the-art methods with a large margin.
	\footnote{Code will be available at https://github.com/ChCh1999/RTPB}


	
	% for each relationship to model the resistance for the model during training. 
	% The RTPB achieves a proper head-tail trade-off to mitigate the bias problem.
	% RTT is motivated by the fact that humans increase muscle strength by making their muscles work against a force. 
	%We assign resistance to the model by using the resistance bias item. RTT can resist the influence of the long tail distribution by 
	% Accordingly, our RTB improves unbiased SGG abilities by forcing network outputs against prior resistant biases (used as the force) for each relationship during training. 
	% The RTB uses prior information from the training set to resist the influence of the long tail distribution on the model and achieve an efficient head-tail trade-off. 
	% Thus we can obtain a balanced model for the scene graph generation. 
	% In addition, to further improve the proposed RTPB by exploring the contextual information, 
\end{abstract}
% Scene graph generation (SGG) is an important task in the computer vision domain, which aims to build a bridge between image data and their semantic meaning and can benefit lots of downstream tasks, like VQA, Image Caption, and Image Retrieve. However, this task suffers from the long-tail distribution of the training set which causes the bias of training. To address this issue, we we propose a novel plug-and-play \textbf{R}esis\textbf{t}ance \textbf{B}ias module (\textit{RTB}) for unbiased scene graph generation. RTB uses prior information from the training set to resist the influence of the long tail distribution on the model and achieve an efficient head-tail trade-off. 
%% Thus we can obtain a balanced model for the scene graph generation. 
%To improve the model performance, we also designed a contextual encoding model Dual Transformer Encoder (DTrans) which encodes the contextual information with transformer. We evaluated our method on all three SGG sub-tasks. Extensive experiments show that the RTB module can achieve significant improvements on the most previous model, which demonstrates the effectiveness of our method. And the combination of DTrans and PB achieves state-of-the-art in mean recall metric.